\documentclass[11pt]{article}
\usepackage[T1]{fontenc}
\usepackage{lmodern,microtype,amsmath,amssymb,amsthm,mathtools,booktabs}
\usepackage[margin=1.08in]{geometry}
\usepackage{xcolor}
\usepackage[colorlinks=true,linkcolor=blue!55!black,citecolor=blue!55!black,urlcolor=blue!55!black]{hyperref}
\hypersetup{pdftitle={Locally resilient rainbow routing with small prime multipliers},pdfauthor={Research note for ProveIt}}
\usepackage{enumitem}
\setlist{nosep}
\newtheorem{theorem}{Theorem}[section]
\newtheorem{lemma}[theorem]{Lemma}
\newtheorem{proposition}[theorem]{Proposition}
\newtheorem{corollary}[theorem]{Corollary}
\theoremstyle{remark}\newtheorem{remark}[theorem]{Remark}
\newcommand{\G}{\mathcal G}
\newcommand{\Pp}{\mathcal P}
\newcommand{\F}{\mathcal F}
\newcommand{\abs}[1]{\lvert#1\rvert}
\newcommand{\norm}[1]{\lVert#1\rVert}
\newcommand{\Z}{\mathbb Z}
\newcommand{\C}{\mathbb C}
\newcommand{\ee}{\varepsilon}
\newcommand{\eps}{\varepsilon}
\newcommand{\ph}{\varphi}
\DeclareMathOperator{\dist}{dist}
\DeclareMathOperator{\Ree}{Re}
\title{Locally resilient rainbow routing\\with small prime multipliers}
\author{Research note for ProveIt}
\date{27 November 2026}
\begin{document}
\maketitle
\begin{abstract}
Assume the Dirichlet generalized Riemann hypothesis in the precise form exported by the vendored Lean development. Let $q$ be a large integer and take the directed graph on $(\Z/q\Z)^\times$ with one labelled arc $u\to up$ for each prime $p\le(\log q)^A$ coprime to $q$. Fix $A>2$ and $0<\delta<1/2$. Even when an adversary independently removes arcs at each residue, retaining at least a $(1/2+\delta)$ fraction of both incoming and outgoing arcs everywhere, every pair of residues is joined by a directed path using distinct primes, of length $O_A(\log q/\log\log q)$. More strongly, every collection of up to $c_{A,\delta}(\log q)^{A-1}$ prescribed disjoint terminal pairs (for some $c_{A,\delta}>0$) can be routed with pairwise vertex-disjoint paths and with no prime used twice anywhere. This order is optimal even without deletions. At cutoff $C_\delta\log^2q$ we obtain paths of length $O_\delta(\log q)$ and capacity $\Omega_\delta(\log q/\log\log q)$. The threshold $1/2$ is asymptotically sharp. The proof is a finite spectral routing lemma combining a two-way random partition of labels, directed expander mixing, and greedy path packing. We give explicit finite inequalities, proofs and a reproducible certificate-generating program. No independence of zeros or nonabelian GRH is used.
\end{abstract}

\section{Statement and significance}
All logarithms are natural. Let $q\ge3$, $L=\log q$, $G=(\Z/q\Z)^\times$, $n=\ph(q)$, and
\[
 \Pp=\{p\le y:p\text{ prime},\ p\nmid q\},\qquad d=\abs{\Pp}.
\]
We always take $y<q$, as is automatic eventually for the cutoffs considered. Thus the prime labels give distinct elements of $G$. Write $\G(q,y)$ for the directed graph with the arcs $x\xrightarrow{p}xp$ for $x\in G$ and $p\in\Pp$, with products interpreted modulo $q$. Each vertex has exactly $d$ arcs in and $d$ out. An admissible graph $H$ is any spanning subgraph, with no requirement that its deletions be invariant under multiplication, for which
\begin{equation}\label{eq:min}
 d_H^+(x),d_H^-(x)\ \ge\ (1/2+\delta)d\qquad (x\in G).
\end{equation}
A rainbow path means a directed simple path whose edge labels are all distinct. In a rainbow linkage we require all vertices of different paths to be distinct and all edge labels across different paths to be distinct as well.

\begin{theorem}[Prime routing and packing]\label{thm:main}
Suppose Dirichlet GRH holds. Fix $A>2$ and $0<\delta<1/2$. Uniformly for all sufficiently large integers $q$, take $y=L^A$. Every admissible $H$ has the following properties.
\begin{enumerate}[label=(\roman*)]
\item Every ordered pair of distinct vertices can be joined by a rainbow path of length at most
\begin{equation}\label{eq:length}
 \left(\frac{2}{A-2}+o(1)\right)\frac{L}{\log L}.
\end{equation}
\item For any $r$ disjoint ordered terminal pairs $(s_i,t_i)$, a rainbow linkage with the same length bound exists provided
\begin{equation}\label{eq:capacity}
 r\le\left(\frac{\delta(A-2)}{64A}-o(1)\right)L^{A-1}.
\end{equation}
Here the vanishing error terms can be chosen independently of $H$ and the pairs; the second assertion means that any fixed coefficient strictly less than the displayed positive coefficient works eventually.
\item If $\mathcal R(q,y,\delta)$ denotes the largest $r\le n/2$ for which every admissible $H$ and every choice of $r$ disjoint ordered pairs admit a rainbow linkage (without a length restriction), then
\begin{equation}\label{eq:theta}
 \mathcal R(q,L^A,\delta)=\Theta_{A,\delta}(L^{A-1}).
\end{equation}
In fact $\mathcal R(q,L^A,\delta)\le(1+o(1))L^{A-1}$, by an obstruction already present in the full graph.
\end{enumerate}
\end{theorem}

\begin{corollary}[At the quadratic-logarithm cutoff]\label{cor:critical}
For each $0<\delta<1/2$ there exist $C_\delta,c_\delta,K_\delta>0$ such that, under Dirichlet GRH, $y=C_\delta L^2$ suffices for a rainbow path of length at most $K_\delta L$ between any two vertices of any admissible $H$. There is a rainbow linkage for any $r\le c_\delta L/\log L$ disjoint pairs, with every path satisfying the same length bound.
\end{corollary}

In arithmetic terms, a route $s\xrightarrow{p_1}\cdots\xrightarrow{p_k}t$ supplies a squarefree $y$-smooth positive integer $P=\prod_jp_j$ with $sP\equiv t\pmod q$. But a restriction may depend on the entire current residue: $p_j$ must be allowed at $s\prod_{i<j}p_i$. Thus ordinary Fourier inversion on a single retained subset of primes does not address this problem. For a linkage the integers $P_i$ are pairwise coprime, and no intermediate residue of one route is used by another. At cutoff $L^A$, each $P_i$ is at most $q^{2A/(A-2)+o(1)}$ by \eqref{eq:length}; the bound $q^{O_\delta(\log\log q)}$ follows at the critical cutoff. These are existence bounds, not least-integer estimates.

\section{Source theorem, analytic input and novelty boundary}
\label{sec:source}
\paragraph{Exact Lean interface.} The requested vendor directory \cite{vendor} contains the readable rewrite of the Alp\"oge--Claude proof, identified by upstream commit
\begin{center}\small\texttt{72fead102f663d5c7a407964a13dc0cf9d977207}.
\end{center}
The types in \texttt{src/Statements.lean} are expressed in Mathlib's language. The exported theorem \texttt{LFunction\_ne\_zero\_of\_half\_lt\_re} says that for all positive moduli and all complex Dirichlet characters, primitive or not, $L(s,\chi)\ne0$ for $\Ree s>1/2$, except at the pole of a principal character. The companion theorem \texttt{LFunction\_eq\_zero\_re} places all zeros in $0<\Ree s<1$ on the critical line. The source chain examined here is
\[
\begin{gathered}
 \texttt{New.Reduction.OneSidedBound}\quad\longrightarrow\quad
 \texttt{New.grh},\\
 \texttt{New.Main}\quad\longrightarrow\quad
 \texttt{Main.lean}.
\end{gathered}
\]
The final proof of \texttt{New.grh} induces a character from $q$ to $18480q$ and invokes the one-sided prime-remainder criterion. The reflected zero argument in \texttt{New.Main} then converts the half-plane theorem into the critical-strip formulation. We use the exported zero statement, not the internal numerical constants or the proof's auxiliary automorphic machinery. In particular we do not apply GRH to general Hecke or Artin $L$-functions.

The vendor README reports the three axioms \texttt{propext}, \texttt{Classical.choice}, \texttt{Quot.sound} and a successful independent-kernel check. The independent verification reports of Sutherland and Rubinstein \cite{sutherland,rubinstein} are consistent with the machine-check status of the source. As of the date of this note, machine verification and formal acceptance by Clay or journal peer review are distinct; Clay lists RH as active \cite{clay}. We take the pinned Dirichlet-GRH statement as input, as the merged article does. We did not rebuild the vendor package (Lean v4.35.0 and its separate Mathlib pin) or kernel-check our new proof.

\paragraph{The merged article.} The requested article \cite{article} is a consolidation of 131 reports (723 pages, 914 results according to its README). The introduction separates unconditional statements, deductions using Dirichlet GRH, and reductions to open estimates. In \texttt{sections/dir.tex}, Lemma \texttt{dir:lem:bias} is exactly the small-character-bias estimate needed below. The quotient theorem, its squarefree-representative corollary and the short-products theorem prove powerful global retained-pool and squarefree-product results. Two later theorems treat simultaneous factorization with quota constraints. Their exact labels are recorded in the accompanying source audit. None of these gives a path with an arbitrary vertex-dependent admissibility relation. Here we keep the small-bias engine but replace fixed-cardinality coefficient estimates by directed set expansion in a non-Cayley subgraph. This distinction is essential: even if every vertex has many outgoing labels, the intersection of the outgoing sets over all vertices may be empty.

\paragraph{Related literature.} Jao, Miller and Venkatesan \cite{jmv} proved GRH-based small-prime Cayley expansion, including for $G=(\Z/q\Z)^\times$. Their spectral input is not claimed as new here. The contribution proposed in this note is the explicit finite local-resilience/rainbow-linkage lemma and the resulting optimal-order routing capacity, with a global no-repeated-prime constraint even across paths. A keyword scan of all 19 substantive article sections, followed by inspection of the relevant product and factorization theorems, found no matching formulation. A focused literature search found the Cayley expansion precedent but did not establish that no other graph-theoretic result implies ours. Novelty is therefore a research claim relative to the supplied article, not a certified priority claim.

\subsection{Uniform character estimates}
Define the (normalized) character bias
\[
 \beta=\max_{1\ne\chi\in\widehat G}\abs{\frac1d\sum_{p\in\Pp}\chi(p)}.
\]
Here and below characters of $G$ are extended by zero on nonunits as Dirichlet characters. We use the following known consequences of Dirichlet GRH; they are restated with the details relevant to uniformity.

\begin{lemma}[Analytic input]\label{lem:bias}
Under Dirichlet GRH, uniformly for all sufficiently large $q$:
\begin{enumerate}[label=(\alph*)]
\item For $y=L^A$, $A>2$ fixed, $d\sim L^A/(A\log L)$ and $\beta\le L^{1-A/2+o(1)}$.
\item For any $b>0$ one can choose fixed $C>0$ so that for $y=CL^2$, $d\sim CL^2/(2\log L)$ and $\beta\le b$.
\end{enumerate}
\end{lemma}
\begin{proof}
Use the standard GRH prime-sum estimate for nonprincipal Dirichlet characters (\cite[Ch.~13, Theorem 13.7 and (13.21)]{mv}; see also \cite{jmv} and the article's \texttt{dir:lem:bias}):
\begin{equation}\label{eq:sharp}
 \abs{\sum_{p\le x}\chi(p)}\ll \sqrt{x}\log(qx)\quad(x\ge2).
\end{equation}
The implied constant is absolute. Imprimitive characters can be handled by the primitive conductor and omission of the primes dividing the larger modulus, at cost $O(\log q)$ for this sum. The prime number theorem and $\omega(q)\le L/\log 2$ give $d\sim y/\log y$ in both regimes. Dividing \eqref{eq:sharp} by $d$ proves (a).

For (b) a sharp cutoff would unnecessarily lose a factor of $\log y$. Instead use the standard integrated explicit formula \cite[Ch.~13, (13.28)]{mv}, in the form
\begin{equation}\label{eq:integrated}
 S_\chi(t):=\sum_{m\le t}(t-m)\Lambda(m)\chi(m)\ll t^{3/2}\log(2q),\quad t\ge 2.
\end{equation}
The nontrivial zero terms have factors $t^{1+\rho}/(\rho(1+\rho))$; on GRH their absolute sum is $O(t^{3/2}\log(2q))$ by the standard zero-count bound. The trivial-zero and elementary terms satisfy the same bound. For imprimitive $\chi$ the additional omitted prime powers contribute at most $t\omega(q)\log t=O(t^{3/2}\log(2q))$. The integrated formula is classical, and is the analytic input explicitly used in the merged article.

Fix $0<\eta<1/2$ and set
\[
 w(m)=\frac{(y-m)_+-((1-\eta)y-m)_+}{\eta y}.
\]
It is one up to $(1-\eta)y$, then tapers linearly to zero. Subtraction in \eqref{eq:integrated} gives\[\abs{\sum_mw(m)\Lambda(m)\chi(m)}\ll\sqrt y L/\eta.\] Prime powers of exponent at least two contribute $O(\sqrt y\log^2y)$, by counting powers and bounding each weight by $\log y$. Restoring the final interval costs at most $(\eta+o(1))y$, by the prime number theorem for $\vartheta$. Replacing each $\log p$ by $\log y$ costs at most $\pi(y)\log y-\vartheta(y)=o(y)$. All these latter estimates are independent of $\chi$. Thus, using $d\log y\sim y$, for $y=CL^2$,
\[
 \beta\le \eta+O(1/(\eta\sqrt C))+o(1).
\]
Choose $\eta$ small, then $C$ large, then $q$ large. This proves (b).
\end{proof}

\section{A finite routing theorem for abelian Cayley graphs}
We now separate the analytic input from the genuinely combinatorial argument. Let $G$ be any finite abelian group of order $n\ge2$. Let $P\subseteq G$ have $d\ge1$ distinct labelled elements, and let $H$ be any spanning subgraph of the arcs $x\to xp$, satisfying \eqref{eq:min}. Define $\beta$ as above, with $P$ replacing $\Pp$. The assumption of distinct elements ensures at most one arc from a vertex to a given other vertex; the argument allows a loop if the identity is a label.

\begin{theorem}[Finite local-resilience lemma]\label{thm:finite}
Let $0<\delta<1/2$ and put
\begin{align*}
 t&=\sqrt{2d\log(16n)}, & \tau&=t/d,\\
 \ee&=2\beta+4\tau, & B&=\left\lfloor\frac{\delta d}{16}\right\rfloor,\\
 R&=\left(\frac{\delta}{2\ee}\right)^2,&
 k&=\left\lfloor\frac{\log(n/2)}{\log R}\right\rfloor+1.
\end{align*}
Assume $\tau\le\delta/16$, $\ee<\delta/2$ (so $R>1$), and $k\le B$. There is a single partition $P=P_+\sqcup P_-$, depending only on $H$, with the following robust property: for every set $F\subseteq P$ of forbidden labels and $W\subseteq G$ of forbidden vertices satisfying
\begin{equation}\label{eq:budget}
 \abs F+\abs W\le B,
\end{equation}
and all $s,t_0\in G\setminus W$, there is a rainbow path from $s$ to $t_0$ of length at most $2k$ avoiding both $F$ and $W$. Before erasing cycles it is a concatenation of a walk of length at most $k$ from $P_+$ and one of length at most $k$ from $P_-$. The same partition works for every choice of $F,W,s,t_0$.

In particular every collection of $r$ disjoint ordered terminal pairs admits a rainbow linkage with paths of length at most $2k$ if
\begin{equation}\label{eq:packing}
 (4k+2)r\le B.
\end{equation}
\end{theorem}

\subsection{Mixing for directed arcs}
For $Q\subseteq P$, write $E_Q(U,V)$ for the number of full-Cayley arcs from $U$ to $V$ with label in $Q$. Set $d_Q=\abs Q$ and $\beta_Q=\max_{\chi\ne1}\abs{d_Q^{-1}\sum_{p\in Q}\chi(p)}$ if $Q\ne\varnothing$.
\begin{lemma}[Directed expander mixing]\label{lem:mix}
For all $U,V\subseteq G$,
\begin{equation}\label{eq:mix}
 \left|\frac{E_Q(U,V)}{d_Q}-\frac{\abs U\abs V}{n}\right|
 \le\beta_Q\sqrt{\abs U\abs V}.
\end{equation}
The same inequality holds with incoming and outgoing interchanged.
\end{lemma}
\begin{proof}
The operator $T_Qf(x)=d_Q^{-1}\sum_{p\in Q}f(xp)$ on $\C^G$ is diagonal in the character basis, with eigenvalue $d_Q^{-1}\sum_{p\in Q}\chi(p)$ on $\chi$. Its operator norm on the orthogonal complement of constants is at most $\beta_Q$, even though its eigenvalues need not be real. In the counting inner product, subtract the constant components from $1_U$ and $1_V$ and apply Cauchy--Schwarz to $\langle 1_U,T_Q1_V\rangle$. The norms of the centred indicators are at most $\sqrt{\abs U}$ and $\sqrt{\abs V}$. The adjoint (reversed arcs) has conjugate eigenvalues and the same bound.
\end{proof}

\subsection{A partition adapted to all local degrees}
\begin{lemma}[Simultaneous split]\label{lem:split}
Under $\tau\le\delta/16$ there is a partition $P=P_+\sqcup P_-$ such that, for each sign $j$,
\begin{align}
 \frac{d-t}{2}\le d_j:=\abs{P_j}&\le\frac{d+t}{2},\label{eq:size}\\
 \beta_{P_j}&\le\ee,\label{eq:poolbias}\\
 d_{H,j}^\pm(x)&\ge(1/2+7\delta/8)d_j\quad(x\in G).\label{eq:splitdeg}
\end{align}
Here $d_{H,j}^\pm(x)$ is the outgoing or incoming degree in $H$ restricted to labels in $P_j$; the sign of the pool is independent of the degree direction.
\end{lemma}
\begin{proof}
For every label choose an independent random sign $\sigma_p\in\{-1,1\}$, defining $P_+$ and $P_-$. For each of the $2n$ sets of incoming or outgoing labels of $H$, require that the absolute signed sum of its indicator be at most $t$. Also, for each of the $n$ characters require that the signed sums of its real and imaginary parts have absolute value at most $t$. Each is a sum $\sum_pa_p\sigma_p$ with $a_p\in[-1,1]$. For completeness, $\mathbb E\exp(u a_p\sigma_p)=\cosh(ua_p)\le\exp(u^2a_p^2/2)$, so optimizing Markov's inequality and applying it to both signs gives a failure probability at most $2\exp(-t^2/(2d))=1/(8n)$. There are at most $4n$ requirements. The union bound shows that a partition satisfying all of them exists (indeed with probability at least $1/2$).

The principal character gives $\abs{\sum_p\sigma_p}\le t$, proving \eqref{eq:size}. For each nonprincipal character the full sum has modulus at most $\beta d$ and the signed sum modulus at most $\sqrt2t$. Consequently
\[
 \beta_{P_j}\le\frac{\beta+\sqrt2\tau}{1-\tau}\le2\beta+4\tau=\ee,
\]
as $\tau\le1/2$. Finally each local retained-label set has at least $(1/2+\delta)d$ labels. In either pool its size is at least $((1/2+\delta)d-t)/2$. Divide by $d_j\le(d+t)/2$ and use
\[
 \frac{1/2+\delta-\tau}{1+\tau}
 \ge1/2+\delta-2\tau\ge1/2+7\delta/8.
\]
This proves \eqref{eq:splitdeg} simultaneously for both directions at every vertex.
\end{proof}

\subsection{Rainbow reachability and growth}
\begin{proof}[Proof of Theorem~\ref{thm:finite}]
Fix the partition from Lemma~\ref{lem:split}, and a choice of forbidden sets satisfying \eqref{eq:budget}. Let $s$ be the source. We describe an exploration restricted to $P_+$; a vertex is stored once, together with a witness directed walk from $s$ whose labels are distinct and avoid $F$, and all of whose vertices avoid $W$. Let $U_0=\{s\}$. Given $U_i$ with stored witnesses of length at most $i$, for every $x\in U_i$ try every outgoing arc in $H$ with label in $P_+$ not in $F$ or the witness's labels, and with endpoint not in $W$. Store one new witness for every new endpoint. Denote by $U_{i+1}$ the resulting enlarged set, keeping all old witnesses. Concatenation preserves the required label condition, even if the walk revisits a vertex. One witness per endpoint suffices; this is not the full set of all possible rainbow walks.

For every $i<k$, a stored witness eliminates at most $i$ labels. The global forbidden labels eliminate at most $\abs F$ arcs from each vertex. Because the generator elements are distinct, excluding vertices in $W$ eliminates at most $\abs W$ arcs from each vertex. Thus the total loss from the minimum in \eqref{eq:splitdeg} is at most $B+i\le2B$. Since $d_+\ge d/3$ by \eqref{eq:size},
\begin{equation}\label{eq:remaining}
 (1/2+7\delta/8)d_+-2B\ge(1/2+\delta/2)d_+.
\end{equation}
All these eligible arcs from $U_i$ end in $U_{i+1}$. They are distinct when counted by their tail and label. Apply \eqref{eq:mix} to get
\[
 (1/2+\delta/2)\abs{U_i}
 \le \frac{\abs{U_i}\abs{U_{i+1}}}{n}
     +\ee\sqrt{\abs{U_i}\abs{U_{i+1}}}.
\]
If $\abs{U_{i+1}}\le n/2$, divide by $\sqrt{\abs{U_i}}$ and square to obtain
\begin{equation}\label{eq:grow}
 \abs{U_{i+1}}\ge R\abs{U_i}.
\end{equation}
If all sets through $U_k$ had size at most $n/2$, then $\abs{U_k}\ge R^k>n/2$ by the definition of $k$, a contradiction. Hence $\abs{U_k}>n/2$ (we can stop once this occurs).

Now run the identical exploration backwards from $t_0$, using incoming arcs whose labels lie in $P_-$; a witness is a directed walk from the current vertex to $t_0$. Excluding $W$ eliminates at most $\abs W$ incoming arcs per vertex since $x p^{-1}$ are distinct predecessors. The same argument gives a set $V$ of size greater than $n/2$ of vertices with reverse witnesses of length at most $k$. There is a vertex $z\in U_k\cap V$. The two witnesses concatenate at $z$ into a walk of length at most $2k$ from $s$ to $t_0$. The labels in the two halves are disjoint because the pools are disjoint; within each half the labels are distinct. Delete any closed subwalk between repeated vertices. This preserves endpoints, allowed arcs and the no-repeated-label property, and terminates with a simple path. All arcs still avoid $F$ and all vertices avoid $W$. When $s=t_0$, the empty path also works.

For the linkage, route the pairs sequentially using this same partition. When routing the $i$th pair, let $F$ be the union of all earlier labels, and let $W$ consist of the earlier internal vertices and all other terminal vertices. Earlier paths avoided the endpoints of all other pairs, so the current terminals are not in $W$. We have
\[
 \abs F\le2k(i-1),\qquad
 \abs W\le2r-2+2k(i-1).
\]
Their sum is at most $(4k+2)r$, hence is within the budget by \eqref{eq:packing}. The one-path result applies at every stage. Induction establishes the pairwise vertex-disjoint and globally label-disjoint linkage.
\end{proof}

\begin{remark}[Why two pools?]
Two independently constructed rainbow walks may use the same prime even when their endpoints meet. Splitting the entire label pool once, while retaining the local degree margin \emph{in both parts}, removes this obstruction. An arbitrary partition, or a partition chosen based only on the global Fourier coefficients, would not suffice for the proof.
\end{remark}

\section{Asymptotics, sharpness and limitations}
\subsection{Upper length and lower packing capacity}
\begin{proof}[Proof of the existence assertions in Theorem~\ref{thm:main}]
With $L=\log q$, $\log n=(1+o(1))L$ uniformly in $q$. Indeed, if $q$ has $m$ distinct prime divisors, then $q\ge(m+1)!$, so $m=O(L/\log L)$. Since $n/q=\prod_{p\mid q}(1-1/p)\ge2^{-m}$, the assertion follows.

Lemma~\ref{lem:bias}(a) gives $d\sim L^A/(A\log L)$ and $\beta\le L^{1-A/2+o(1)}$. Therefore
\[
 \tau=O\left(L^{(1-A)/2}\sqrt{\log L}\right),\quad
 \ee\le L^{1-A/2+o(1)},\quad
 \log R\ge (A-2-o(1))\log L.
\]
All the finite-lemma hypotheses hold eventually: $\tau\to0$, $\ee\to0$, and
\[
 k\le\left(\frac{1}{A-2}+o(1)\right)\frac{L}{\log L}=o(d),
 \qquad B=\left\lfloor\frac{\delta d}{16}\right\rfloor.
\]
The two-way path bound is \eqref{eq:length}. From the finite packing condition,
\[
 \left\lfloor\frac{B}{4k+2}\right\rfloor
 \ge\left(\frac{\delta(A-2)}{64A}-o(1)\right)L^{A-1},
\]
where floors and additive constants are absorbed by the vanishing error, proving \eqref{eq:capacity}.
\end{proof}

\begin{proof}[Proof of Corollary~\ref{cor:critical}]
Choose $C_\delta$ in Lemma~\ref{lem:bias}(b) with $b=\delta/64$. Then $\tau=O_\delta(\sqrt{\log L/L})\to0$, so for all sufficiently large $q$ we have $\tau\le\delta/16$ and $\ee\le\delta/8$. Thus $R\ge16$ and $k\le L/\log 16+1$. Since $d\sim C_\delta L^2/(2\log L)$, $k\le B$ eventually and $B/(4k+2)\gg_\delta L/\log L$. Apply Theorem~\ref{thm:finite}. No claim of optimal routing capacity at this critical cutoff is made.
\end{proof}

\subsection{A counting obstruction to larger linkages}
\begin{proposition}[Universal capacity upper bound]\label{prop:upper}
Let $G$ have $n\ge4$ vertices and $d\ge2$ distinct generators. Set
\[
 h=1+\left\lfloor\frac{\log(n/4)}{\log d}\right\rfloor.
\]
If $r=\lfloor d/h\rfloor+1$ satisfies $2r\le n/2$, there exist $r$ disjoint ordered pairs with no globally label-disjoint linkage, even in the full directed Cayley graph.
\end{proposition}
\begin{proof}
From any source, at most $1+d+\cdots+d^{h-1}\le 2d^{h-1}\le n/2$ vertices are reachable in fewer than $h$ steps, counting all labelled walks, including repeated labels. Choose sources and targets greedily: at each step pick an unused source, then pick an unused target outside this short-distance ball. There is such a target since at most $2r-1<n/2$ previously selected or current terminals are excluded. Every path for each of these pairs uses at least $h$ labels. Globally distinct labels would require at least $rh>d$, which is impossible.
\end{proof}

In the regime $y=L^A$ we have $h\sim L/(A\log L)$ and $d/h\sim L^{A-1}$. The bound on $2r$ in the proposition holds for large $q$ since $n=q^{1-o(1)}$. A failing collection can be extended to any larger size up to $\lfloor n/2\rfloor$ by adding disjoint pairs, so the capacity is at most $\lfloor d/h\rfloor$. This proves the upper bound in Theorem~\ref{thm:main}(iii). The same counting argument shows that the worst-case one-pair length is $\Omega_A(L/\log L)$, even in the full graph. Thus the order of both the length and the capacity is best possible in the supercritical regime. The constants need not be optimal. At the critical cutoff the capacity upper bound obtained this way is only $O_\delta(L)$; our lower bound leaves a factor of order $\log L$.

\subsection{The sharp local-deletion threshold}
\begin{proposition}[Obstructions]\label{prop:sharp}
For every $q>2$ there is a nontrivial quadratic character of $G=(\Z/q\Z)^\times$. In any regime with $\beta=o(1)$, one can disconnect the directed prime Cayley graph while retaining $(1/2+o(1))d$ outgoing and incoming arcs at every vertex. Also, retaining nearly all outgoing arcs alone, without an incoming requirement, does not guarantee reachability, and conversely.
\end{proposition}
\begin{proof}
The order $n$ is even (multiplication by $-1$ pairs distinct units). Every finite abelian group of even order admits a quotient of order two, hence a character $\chi:G\to\{\pm1\}$ onto that group. Keep precisely the arcs with $\chi(p)=1$. They never leave the coset of $\ker\chi$ containing the tail. Each vertex has the same incoming and outgoing degree, namely
\[
 d_+=\frac{d+\sum_{p\in\Pp}\chi(p)}2=\frac d2+O(\beta d).
\]
This proves the first obstruction. At cutoff $CL^2$, Lemma~\ref{lem:bias}(b) allows this retained fraction to be made arbitrarily close to $1/2$ by increasing the fixed $C$ before taking $q$ large; it does not assert $\beta=o(1)$ for each fixed $C$. For the one-sided-degree obstruction, delete every incoming arc of a designated vertex. Each other vertex loses at most one outgoing arc (distinct generators), but the designated vertex cannot be reached from any other. Deleting outgoing arcs instead proves the converse. These examples explain both the threshold and the need for the two directions in \eqref{eq:min}.
\end{proof}

\section{Reproducibility and formalization}
The accompanying \texttt{routing.py} constructs locally damaged graphs and implements the two-sided witness exploration. It uses Python's standard library. In the three included examples, the sets of retained outgoing prime labels vary from vertex to vertex. The file \texttt{examples.json} contains all terminal pairs, partitions, routes, and their intermediate residues. The verifier reconstructs the edge-deletion instance from its random seed and checks every directed multiplication, membership in the retained graph, route length, and pairwise disjointness of vertices and prime labels. Its checks are exact integer computations; they do not assume GRH or estimate eigenvalues numerically. To reproduce, run
\begin{center}\small\texttt{python3 routing.py --output examples.json}\\
\texttt{python3 routing.py --verify examples.json}.
\end{center}
The exact results were:
\begin{center}
\begin{tabular}{rrrrrr}
\toprule
$q$ & $n$ & $d$ & Min. in/out degrees & Routes & Max. length\\
\midrule
1009 & 1008 & 109 & 82 / 81 & 10 & 2\\
1155 & 480 & 105 & 78 / 76 & 10 & 2\\
5003 & 5002 & 154 & 114 / 115 & 10 & 3\\
\bottomrule
\end{tabular}
\end{center}
The small examples illustrate the construction; the conservative finite inequalities of Theorem~\ref{thm:finite} need not hold at these sizes. Their validity is not inferred from the experiments.

For an explicit finite graph satisfying the theorem's inequalities one may randomly try partitions until the linkage succeeds, verifying each output. A good partition occurs with probability at least $1/2$ by Lemma~\ref{lem:split}, and any such partition works for all requested pairs within budget. The exploration needs at most $k$ rounds per direction, each inspecting at most $nd$ arcs; the path certificate itself has at most $2kr$ arcs. The demonstration program additionally tests local split degrees and uses a bounded retry count. This is polynomial in the explicitly represented graph size, not a claim of polynomial time in the bit length of $q$.

A Lean development could separate the analytic and finite interfaces: use \texttt{Main} from the vendored Lake package to instantiate Dirichlet GRH; derive \eqref{eq:sharp} and \eqref{eq:integrated} from a formal explicit formula; formalize character orthogonality on a finite abelian group and the Hoeffding union bound; encode witness walks as finite lists of arcs; prove the counting and greedy packing lemmas on finite sets. The analytic estimates are classical imports in the present paper, not Lean results verified here. See \texttt{FORMALIZATION.md} for a dependency plan and \texttt{SOURCE\_AUDIT.md} for the inspected blob identifiers and scope.

\section{Conclusion and further questions}
We have moved from globally retained sets of prime factors to an arbitrary local directed transition relation. A two-pool spectral argument makes the global prime-disjointness constraint compatible with strong local damage; its capacity is matched in order by the finite number of available labels. Two natural questions remain: can one match the capacity upper bound $O(L)$ at the cutoff $CL^2$, and can one improve the constant $2/(A-2)$ in the length estimate while retaining the same adversarial threshold? The present proof does not resolve these questions or any non-Dirichlet Riemann hypothesis.

\begin{thebibliography}{9}
\bibitem{vendor} L. Sernau and collaborators, \emph{The Generalized Riemann Hypothesis for Dirichlet L-functions, in Lean 4}, readable edition, 25 Nov. 2026, upstream commit \texttt{72fead10}. \href{https://github.com/VladimirReshetnikov/ProveIt/tree/main/Analysis/GeneralizedRiemannHypothesis/vendor/grh-lean}{Vendored directory}; \href{https://github.com/lsernau/grh-lean/tree/72fead102f663d5c7a407964a13dc0cf9d977207/grh}{upstream pinned source}.
\bibitem{article} ProveIt Contributors, \emph{Beyond the Riemann hypothesis for Dirichlet L-functions: consequences, transfers and limits}, 26 Nov. 2026, in particular \texttt{sections/src.tex} and \texttt{sections/dir.tex}. \href{https://github.com/VladimirReshetnikov/ProveIt/tree/main/Analysis/GeneralizedRiemannHypothesis/Article}{Article directory}.
\bibitem{jmv} D. Jao, S. D. Miller and R. Venkatesan, Expander graphs based on GRH with an application to elliptic curve cryptography, \emph{J. Number Theory} \textbf{129} (2009), 1491--1504. \href{https://arxiv.org/abs/0811.0647}{arXiv:0811.0647}.
\bibitem{mv} H. L. Montgomery and R. C. Vaughan, \emph{Multiplicative Number Theory I: Classical Theory}, Cambridge University Press, 2007, Chapter 13.
\bibitem{sutherland} A. V. Sutherland, verification files for the shorter Lean proof of the Riemann hypothesis, announced 25 Nov. 2026. \href{https://math.mit.edu/~drew/}{MIT research page}.
\bibitem{rubinstein} M. Rubinstein, Alp\"oge/Claude proof of the Generalized Riemann Hypothesis, including independent verification reports, November 2026. \href{https://www.math.uwaterloo.ca/~mrubinst/alpoge_claude_grh/}{Waterloo research page}.
\bibitem{clay} Clay Mathematics Institute, \href{https://www.claymath.org/millennium/Riemann-Hypothesis/}{Riemann Hypothesis}, status viewed 27 Nov. 2026.
\end{thebibliography}
\end{document}
